\section{Costo Total Estimado}

El costo total del proyecto resulta de la suma de las dos partidas estimadas:

\begin{equation}
    C_{\text{total}} = C_{\text{software}} + C_{\text{equipamiento}}
\end{equation}

La Tabla \ref{tab:costo_total} presenta la estructura de la estimación.

\begin{longtable}{
        P{0.25}
        P{0.52}
        P{0.15}
    }
    \caption{Estructura del costo total estimado}
    \label{tab:costo_total}\\

    \toprule
    \rowcolor{gray!20}
    \textbf{Partida} &
    \textbf{Composición} &
    \textbf{Valor} \\
    \midrule
    \endfirsthead

    \multicolumn{3}{c}{\textit{Continuación de la Tabla
    \ref{tab:costo_total}}} \\
    \toprule
    \rowcolor{gray!20}
    \textbf{Partida} &
    \textbf{Composición} &
    \textbf{Valor} \\
    \midrule
    \endhead

    Costo de software &
    $62$ puntos de función $\times$ $12$ horas por punto $\times$ Bs $30$ por hora &
    Bs 22,320 \\
    \midrule

    Costo de equipamiento &
    Equipo de cómputo con acelerador gráfico y almacenamiento externo &
    Bs 10,850 \\
    \midrule

    \multicolumn{2}{r}{\textbf{Costo total del proyecto}} & \textbf{Bs 33,170} \\*
    \multicolumn{3}{r}{\footnotesize\textit{Fuente: Elaboración propia}} \\*
    \bottomrule
\end{longtable}

Dos observaciones corresponden a esta estructura. La primera es que la partida de
software domina la estimación ---cerca del $67\%$ del total--- porque el proyecto
no requiere adquirir infraestructura de captura ni licencias, de manera que el
valor del desarrollo concentra la mayor parte del costo. La segunda es que la
estimación se sostiene sobre \textbf{parámetros declarados} ---productividad,
tarifa horaria y precios de equipo--- y no sobre supuestos implícitos, lo que
permite recalcularla ante una variación de las condiciones de mercado sin
modificar el análisis técnico.

La valoración monetaria queda así cerrada con los parámetros declarados: Bs
33,170 como costo total, de los cuales Bs 22,320 corresponden al desarrollo y
Bs 10,850 al equipamiento. Desde el punto de vista metodológico, la estimación
establece un tamaño funcional verificable ---$62$ puntos de función--- corroborado
por una medición independiente del código ---$6{,}946$ líneas de producto en $88$
archivos---, que constituye la base trazable de cualquier valoración posterior.
